%%%%%%%%%%%%%%%%%%%%%%%%%%%%%%%%%%%%%%%%%%%%%%%%%%%%%%%%%%%%%%%%%%%%%%%%
%% QUADRO REIVINDICATÓRIO                                              %%
%%                                                                     %%
%% Três regras de forma valem para TODA reivindicação (art. 28):        %%
%%   I   numerada consecutivamente, em algarismos arábicos — o ambiente %%
%%       'quadroreivindicatorio' cuida disso;                           %%
%%   II  conter uma ÚNICA expressão "caracterizado por";                %%
%%   III ser redigida sem interrupção por pontos — um único ponto final.%%
%%                                                                     %%
%% Escreva cada reivindicação após um \reivindicacao. Cuidado com o     %%
%% inciso III: nada de ponto no meio do texto, e portanto nada de       %%
%% abreviaturas com ponto ou de números com ponto decimal (use vírgula).%%
%%                                                                     %%
%% Estrutura de uma reivindicação INDEPENDENTE (art. 30, V e VI):       %%
%%     Preâmbulo + "caracterizado por" + matéria pleiteada              %%
%% Estrutura de uma reivindicação DEPENDENTE (art. 31, I):              %%
%%     Preâmbulo + "de acordo com a reivindicação N" +                  %%
%%     "caracterizado por" + matéria pleiteada                          %%
%%                                                                     %%
%% Art. 29, VIII — as características técnicas devem vir acompanhadas   %%
%% dos sinais de referência das figuras, entre parênteses, quando       %%
%% essencial à clareza. Em modelo de utilidade isso é obrigatório       %%
%% (art. 32, III).                                                      %%
%%                                                                     %%
%% Art. 29, VII e IX — não faça referência ao relatório ou aos desenhos %%
%% ("como descrito na parte ... do relatório") e não inclua trechos     %%
%% explicativos sobre vantagens ou sobre o simples uso do objeto.       %%
%%                                                                     %%
%% Arts. 10 e 18 da LPI — cuidado com a matéria que a lei não considera %%
%% invenção ou não considera patenteável. Não se pleiteia extrato       %%
%% vegetal (art. 10, IX), mas pleiteia-se a composição que o contém;   %%
%% não se pleiteia programa de computador em si (art. 10, V), mas       %%
%% pleiteia-se o produto ou o processo que o utiliza.                   %%
%%                                                                     %%
%% Por fim: tudo o que você pleiteia aqui precisa estar descrito e      %%
%% concretizado no relatório descritivo, sob pena de o pedido ser       %%
%% indeferido por não atender aos arts. 24 e 25 da LPI — mesmo que o    %%
%% objeto seja realmente novo e inventivo.                              %%
%%%%%%%%%%%%%%%%%%%%%%%%%%%%%%%%%%%%%%%%%%%%%%%%%%%%%%%%%%%%%%%%%%%%%%%%

\begin{quadroreivindicatorio}

\ifinvencao

%% ---------------------------------------------------------------------
%% PATENTE DE INVENÇÃO — arts. 29 a 31
%%
%% Art. 30, I — para cada categoria de reivindicação cabe pelo menos uma
%%   reivindicação independente.
%% Art. 29, III — um mesmo quadro pode conter mais de uma categoria
%%   (produto e processo, processo e aparelho, produto e uso...), desde que
%%   ligadas pelo mesmo conceito inventivo (art. 22 da LPI).
%% Art. 30, IV — quando as categorias se interligam, a parte inicial da
%%   reivindicação deve evidenciar a interligação: "Processo para a obtenção
%%   do produto definido na reivindicação...".
%% Art. 30, VII e art. 31, VI — agrupe as dependentes por categoria, para
%%   dar concisão ao quadro.
%% Art. 31, III — não se admite "de acordo com uma ou mais das
%%   reivindicações..." nem "de acordo com as reivindicações anteriores...";
%%   "de acordo com qualquer uma das reivindicações anteriores" é aceito.
%% ---------------------------------------------------------------------

\reivindicacao DISPOSITIVO DE ACIONAMENTO PARA VÁLVULA DE IRRIGAÇÃO,
compreendendo um corpo tubular (2) provido de bocal de entrada e de bocal de
saída, uma membrana elástica (3) alojada no interior do referido corpo tubular
(2) e solidária a uma haste de acionamento (4), e um atuador eletromecânico (6)
acoplado à referida haste de acionamento (4), caracterizado por a haste de
acionamento (4) apresentar em sua porção intermediária um ressalto de encosto
atuante sobre uma mola de retorno (5) pré-carregada e receber, em sua
extremidade oposta à membrana elástica (3), um assento de vedação de dupla face
(9) que veda alternadamente o bocal de saída e um orifício de alívio do corpo
tubular (2), sendo a pré-carga da mola de retorno (5) dimensionada para manter
a haste de acionamento (4) em duas posições extremas estáveis sem alimentação
do atuador eletromecânico (6).

\reivindicacao DISPOSITIVO DE ACIONAMENTO PARA VÁLVULA DE IRRIGAÇÃO, de acordo
com a reivindicação 1, caracterizado por compreender um sensor de pressão (7)
montado no bocal de saída do corpo tubular (2) e uma unidade de controle (8)
conectada ao referido sensor de pressão (7) e ao atuador eletromecânico (6).

\reivindicacao DISPOSITIVO DE ACIONAMENTO PARA VÁLVULA DE IRRIGAÇÃO, de acordo
com a reivindicação 2, caracterizado por a unidade de controle (8) comandar o
atuador eletromecânico (6) por meio de pulsos de corrente com duração inferior
a duzentos milissegundos.

\reivindicacao PROCESSO DE CONTROLE DE VAZÃO, executado pelo dispositivo de
acionamento definido na reivindicação 2, compreendendo as etapas de leitura do
sinal de pressão fornecido pelo sensor de pressão (7) e de comparação do valor
lido com uma faixa de operação previamente definida, caracterizado por
compreender ainda a etapa de aplicação de um pulso de corrente ao atuador
eletromecânico (6) quando o valor lido permanecer fora da faixa de operação por
intervalo superior a um tempo de confirmação.

\reivindicacao PROCESSO DE CONTROLE DE VAZÃO, de acordo com a reivindicação 4,
caracterizado por a faixa de operação ser delimitada por um limite inferior e
um limite superior de pressão\incluido{, ambos ajustáveis pela unidade de
controle (8)}.

%% Exemplo de reivindicação suprimida (art. 57, II). No documento do pedido não
%% imprime nada e não consome número, de modo que as demais fiquem numeradas
%% consecutivamente como exige o art. 28, I; na cópia de comparação sai tachada,
%% exibindo o número que tinha no quadro anterior.
\reivindicacaoremovida{6}{USO DO DISPOSITIVO DE ACIONAMENTO definido na
reivindicação 1, caracterizado por ser em sistema de irrigação por gotejamento.}

\else

%% ---------------------------------------------------------------------
%% PATENTE DE MODELO DE UTILIDADE — arts. 32 a 36
%%
%% Art. 33 — cada pedido deve conter UMA ÚNICA reivindicação independente,
%%   que descreva o modelo definindo integralmente todas as características
%%   de forma ou disposição introduzidas, essenciais à obtenção da melhoria
%%   funcional.
%% Art. 32, I — as reivindicações devem ser preferencialmente iniciadas
%%   pelo título do pedido.
%% Art. 32, III — os sinais de referência das figuras são obrigatórios,
%%   entre parênteses.
%% Art. 34, II — depois de "caracterizado por" vem SOMENTE a nova forma ou
%%   disposição introduzida, com todos os elementos que a constituem e seus
%%   posicionamentos e interconexões em relação ao conjunto.
%% Art. 32, VI — não são aceitas reivindicações de utilização nem trechos
%%   explicativos sobre vantagens ou sobre o uso do objeto.
%% Art. 35 — dependentes só são aceitas quando se referirem a: (I) elemento
%%   complementar de uso opcional que não altere as condições de utilização
%%   e funcionamento; (II) variação de forma ou detalhe de elementos
%%   componentes definidos na primeira reivindicação, sem alterar a unidade
%%   do modelo; ou (III) o objeto em sua forma tridimensional, quando a
%%   configuração final decorra da montagem de estrutura inicial planificada.
%% ---------------------------------------------------------------------

\reivindicacao DISPOSITIVO DE ACIONAMENTO PARA VÁLVULA DE IRRIGAÇÃO,
constituído por um corpo tubular (2) provido de bocal de entrada e de bocal de
saída, por uma membrana elástica (3) alojada no interior do referido corpo
tubular (2) e solidária a uma haste de acionamento (4), e por um atuador
eletromecânico (6) acoplado à referida haste de acionamento (4), caracterizado
por a haste de acionamento (4) apresentar em sua porção intermediária um
ressalto de encosto assentado contra uma mola de retorno (5) pré-carregada
apoiada na face interna do corpo tubular (2), e por receber, em sua extremidade
oposta à membrana elástica (3), um assento de vedação de dupla face (9)
posicionado entre o bocal de saída e um orifício de alívio do corpo tubular
(2).

\reivindicacao DISPOSITIVO DE ACIONAMENTO PARA VÁLVULA DE IRRIGAÇÃO, de acordo
com a reivindicação 1, caracterizado por o ressalto de encosto da haste de
acionamento (4) apresentar formato troncônico.

\fi

\end{quadroreivindicatorio}
